% !TEX root = ../main.tex

\section{INTRODUCTION}
\label{sec:intro}

% =====================================================================
% 26 Oct 2026: plain-language rewrite from the team's Google Doc draft.
% Problems, research questions and objectives follow the same 1-4 order
% as Chapter 3. Hardware (Tesla T4 via Colab) and the stabilization-vs-SNR
% comparison follow Chapter 3. Every acronym is listed in frontmatter.tex.
% =====================================================================

\subsection{Background of the Study}
\label{sec:background}

Giving computers a voice is the goal of text-to-speech (TTS) technology, which turns written text into spoken output while aiming to keep the speech clear, true to the speaker's voice, and natural. One approach used in generative TTS systems is flow matching, which describes a smooth change from random noise into a representation of speech. Flow-matching models learn how to move along a path that connects noise to speech, and this path is often fairly simple or close to a straight line. The learned directions are then followed through a series of calculations to produce speech \parencite{lipman2023flow,liu2023flow}. F5-TTS, developed by \textcite{chen2024f5}, uses flow matching through a Diffusion Transformer and introduces Sway Sampling, which changes the path followed when speech is generated. These developments form a basis for studying ways to generate speech more efficiently, especially because the full computing cost of such a model may not be needed for every input.

Despite these advances, generating speech can still require many repeated calculations by the neural network. This increases computing cost and may slow the response, which matters in interactive applications such as communication aids. Knowledge distillation addresses this challenge by training a smaller or faster student model to imitate a more computationally expensive teacher. \textcite{li2024dmospeech} developed DMOSpeech, a zero-shot TTS system that uses distribution-matching distillation and reduces the generation process from 128 steps to four. Similarly, \textcite{zhu2025zipvoice} introduced ZipVoice, a flow-matching TTS system with a distilled version designed to speed up speech generation. These studies show how distillation can lower the computing needs of speech generation. However, the efficiency of a distilled model does not show that a single setting gives the best balance between computing cost and speech quality for every input.

Building on this direction, \textcite{li2025dmospeech2} developed DMOSpeech~2, a teacher-guided TTS system in which a teacher model carries out the first generation steps and a distilled student model completes the speech. The authors report that this approach improves output diversity compared with using the student alone, while remaining more efficient than running the full teacher. This provides a direct precedent for combining teacher guidance and distilled generation in TTS. However, the handover from the teacher to the student follows a predefined switching time rather than a method that selects the switching time separately for each utterance. As a result, the amount of computing assigned to the teacher stays fixed, even though different inputs may have different needs. Studying whether the handover time can be selected adaptively may therefore offer a way to adjust computing effort to each input while keeping the published generation procedure of the distilled student.

Adaptive cooperation between teacher and student models has also been studied in other generative tasks. \textcite{starodubcev2024your} proposed a method for generating images from text in which a student first creates an image and a quality estimator decides whether the teacher needs to refine it further. This approach shows how teacher computing can be used selectively rather than equally for all outputs. \textcite{mualem2026aria} introduced ARIA, which adjusts how much training effort is given according to the differences between the teacher and student across parts of the input. Although these studies offer relevant examples of adjusting computing effort and resources, they address different tasks or stages of model development. \textcite{starodubcev2024your} focus on improving images during generation, whereas \textcite{mualem2026aria} concern the distillation process itself. The reviewed sources therefore provide related ideas but do not establish a method for selecting the teacher-to-student handover time for individual utterances in flow-matching TTS.

An adaptive handover policy also raises the question of how reliable its decisions are. Choosing a switching time based on signals available during generation does not, by itself, show how often the chosen setting may introduce additional quality failures. Learn-Then-Test (LTT), introduced by \textcite{angelopoulos2025ltt}, is a statistical framework for choosing among a limited set of candidate settings while making sure, at a chosen confidence level, that the risk of the chosen setting does not go beyond a stated limit, as long as the framework's assumptions are met. Applying LTT to teacher-guided TTS requires defining what counts as a quality failure, choosing a reference setting, stating the acceptable risk limit, and using suitable calibration data. In this study, LTT is meant to control the measured rate of additional quality failures, not to guarantee the quality of every individual utterance.

The reviewed literature provides the foundations for this study through flow-matching TTS, distilled speech generation, teacher-guided sampling, adaptive computation, and statistical risk calibration. However, within the sources examined, no single method has been established that combines the selection of the teacher--student handover time for each utterance with statistical calibration of the added risk of quality failures in flow-matching TTS. This gap is limited to the combination of these elements in the literature reviewed and does not mean that no other adaptive methods exist. The study therefore investigates whether an adaptive handover policy can improve how teacher computing is assigned while keeping the added quality failures at a statistically controlled level.

To address this problem, the study proposes a training-free, risk-calibrated adaptive handover method that uses the distilled model of DMOSpeech~2 \parencite{li2025dmospeech2}. The proposed method will use signals available during generation to choose among candidate handover times while keeping the pretrained model weights unchanged. Instead of adjusting the student's restart level as an active variable during generation, the model keeps its standard, predefined default setting. LTT will be used to calibrate the eligible handover settings against a predefined limit for additional quality failures relative to a reference setting, subject to the framework's statistical assumptions. The method will be evaluated using word error rate (WER), speaker similarity (SIM), predicted naturalness (UTMOS), forward-pass counts, also called neural function evaluations (NFEs), and runtime. These measures will be used to determine whether adaptive handover scheduling can reduce computing cost while controlling the measured risk of additional quality failures.

Efficient and responsive speech generation may also be relevant to assistive information and communication technology for autistic users. Research on augmentative and alternative communication (AAC), including speech-generating devices, has looked at their role in supporting communication among autistic individuals \parencite{therrien2025sgd,muharib2024sgd}. A TTS component can provide spoken output for text that a user selects or types, which makes the efficiency of speech generation relevant to communication-support applications. In this study, the assistive-technology context of autism spectrum disorder (ASD) motivates the practical value of developing a more computationally efficient TTS method. The main evaluation, however, focuses on speech quality, computing cost, and the statistical control of additional quality failures. It does not assess clinical outcomes or establish how effective the resulting application is as an AAC intervention. Experiments will use an NVIDIA Tesla T4 graphics processing unit (GPU) through Google Colab, while use on mobile or other devices with limited computing power remains a possible direction for future work.

% =====================================================================
\subsection{Problem Statement}
\label{sec:problem}

Speaking without delay matters most when a person relies on a device to communicate. DMOSpeech~2, developed by \textcite{li2025dmospeech2}, produces speech by letting a teacher model carry out the first generation steps and a distilled student model finish the speech. The switch between the two happens at a predefined time. However, a fixed handover may not give the best balance between speech quality and computing cost for every utterance. Choosing the handover time adaptively also raises questions about how it affects the speech generation process and how reliable its decisions are. These concerns matter for efficient TTS, including its possible use in communication-support applications, such as those intended for autistic learners.

To begin with, the handover changes the way speech is generated. At the handover, the student takes the teacher's estimate of the final result and combines it with fresh noise according to its restart setting \parencite{li2025dmospeech2}. Changing the switching time therefore changes what the student receives. DMOSpeech~2 reports a single teacher-guided setting with results averaged over its test sets \parencite{li2025dmospeech2}, so how different switching times affect speech quality and computing cost, and how these effects differ from one utterance to another, needs to be tested through experiments, particularly in terms of intelligibility, speaker similarity, and predicted naturalness.

Building on this, the handover point itself is a concern. A predefined switching time decides how much computing the teacher does before the student finishes the speech \parencite{li2025dmospeech2}. Because each guided teacher step needs two network evaluations \parencite{ho2022cfg}, a switching time that is later than an utterance needs spends computing that may not be required, while one that is too early may harm the speech. It is still unclear whether choosing among candidate handover times, using signals available while the speech is being generated, can lower computing cost while keeping speech quality across different utterances.

Beyond these effects on quality and cost, there is no measure of reliability for adaptive handover. Choosing the handover adaptively does not, by itself, show how often additional quality failures occur compared with a fixed reference setting. Learn-Then-Test provides a framework for controlling risk statistically when choosing among candidate settings, as long as its assumptions are met \parencite{angelopoulos2025ltt}. Whether it can be applied to adaptive TTS handover must therefore be examined using a defined quality-failure criterion, calibration data kept separate from the evaluation data, and a risk limit set beforehand.

Finally, these technical concerns connect to the use of speech synthesis for neurodivergent learners. Research has examined speech-generating devices as communication-support tools for autistic individuals \parencite{therrien2025sgd,muharib2024sgd}. However, the possible relevance of TTS to these uses does not show that a more computationally efficient method improves communication outcomes. The proposed method must therefore be judged on its technical performance alone: against fixed switching times and simpler rules, under both clean and noisy prompts, and without assuming that it is usable or effective as an assistive intervention.

In view of these concerns, the general problem addressed by this study is whether a training-free, risk-calibrated adaptive handover method can reduce computing cost in teacher-guided flow-matching TTS while keeping the added risk of quality failures, compared with a fixed reference setting, within a limit set beforehand. Specifically, the study seeks to answer the following research questions:
\begin{enumerate}[leftmargin=1.4em]
\item How do different switching times affect speech quality (intelligibility, speaker similarity, and predicted naturalness) and computing cost, and how do these effects vary across utterances, when the student restarts from the teacher's estimate combined with fresh noise at its default restart level?
\item Can choosing among candidate handover times, using signals available during speech generation, lower computing cost while keeping speech quality across different utterances, without adding network evaluations or changing any model weights?
\item Can the Learn-Then-Test framework, using a defined quality-failure criterion, calibration data, and a risk limit set beforehand, control the rate of additional quality failures of adaptive handover relative to a fixed reference setting?
\item How does the calibrated method compare technically with fixed switching times, a rule based only on the prompt's signal-to-noise ratio (SNR), and an offline oracle in speech quality, computing cost, and risk control under clean and noisy prompts, and does the endpoint-stabilization score add value beyond the prompt SNR?
\end{enumerate}

% =====================================================================
\subsection{Objectives of the Study}
\label{sec:objectives}

\subsubsection{General Objective}
The general objective of this study is to develop and evaluate a training-free, risk-calibrated method for choosing the teacher--student handover setting for each utterance in teacher-guided flow-matching TTS, aiming to lower computing cost while keeping speech quality within a risk limit, set beforehand, for additional quality failures, with all model weights kept unchanged.

\subsubsection{Specific Objectives}
Specifically, this study aims to:
\begin{enumerate}[leftmargin=1.4em]
\item \textbf{Describe the handover settings.} Measure how word error rate, speaker similarity, predicted naturalness, forward-pass counts, and running time change across candidate switching times and input utterances, using the published teacher-guided sampling procedure while keeping the student's default restart level.
\item \textbf{Formulate and analyze the adaptive selection method.} Define the choice of handover for each utterance as a problem of minimizing cost under a set limit on risk, develop a training-free selection rule that uses signals available during generation without changing the pretrained model weights or adding network evaluations, and analyze its properties.
\item \textbf{Calibrate the method for reliability.} Develop and apply a calibration procedure based on Learn-Then-Test to keep the risk of additional quality failures, compared with a fixed reference setting, within a limit and confidence level set beforehand, and check the resulting guarantee on held-out evaluation data.
\item \textbf{Compare the decision signals and the quality--cost trade-off.} Compare versions of the method that use the prompt SNR, the endpoint-stabilization score, or both, each calibrated with the same procedure, against fixed switching times and an offline oracle, using word error rate, speaker similarity, predicted naturalness, forward-pass counts, and runtime under clean and noisy prompt conditions; determine whether the stabilization score adds value beyond the prompt SNR; and examine how the computational savings and the reliability guarantee vary with calibration sample size.
\end{enumerate}

% =====================================================================
\subsection{Significance of the Study}
\label{sec:significance}

This study investigates a training-free, risk-calibrated adaptive handover method for teacher-guided flow-matching TTS. Its findings may contribute to more efficient speech synthesis while keeping speech quality within a defined risk limit. The results are expected to benefit the following:
\begin{description}[leftmargin=0pt,itemsep=2pt,font=\normalfont\itshape]
\item[TTS researchers.] The study may provide insights into how the choice of handover affects speech quality and computing cost, supporting further research on adaptive TTS inference.
\item[TTS developers.] The proposed method may offer a way to choose the handover setting for each utterance without retraining the pretrained models, potentially reducing computing cost while controlling additional quality failures.
\item[Researchers in reliable machine learning.] The application of Learn-Then-Test may show how statistical risk control can be used to choose inference settings adaptively under stated calibration assumptions.
\item[Assistive-technology and AAC developers.] The findings may inform future work on efficient TTS components for communication-support technologies, including augmentative and alternative communication.
\item[Autistic and neurodivergent learners.] The study may indirectly support future assistive-technology research by examining the technical performance of speech synthesis. However, it does not directly evaluate usability or communication outcomes for these learners.
\item[Future researchers.] The methodology and findings may serve as a reference for further studies on adaptive handover, risk-calibrated inference, and efficient TTS across different datasets, noise conditions, and hardware.
\end{description}

% =====================================================================
\subsection{Scope and Limitations}
\label{sec:scope}

This study focuses on developing and evaluating a training-free, risk-calibrated method for choosing the teacher--student handover setting for each utterance in teacher-guided flow-matching TTS. It covers the choice of switching times, the signals available during generation that guide the decision, the selection rule, statistical calibration, and objective evaluation of speech quality and computing cost. The switching time is the only handover setting that the study varies.

All pretrained model weights, including those of the teacher, the student, the duration predictor, and the vocoder, are kept unchanged. Duration prediction follows the published procedure and is not part of the proposed contribution. The student's restart level is likewise kept at the released default for every switching time, so that differences between candidate switching times reflect only how long the teacher runs. The evaluation uses English zero-shot prompts from the Seed-TTS test set, which draws on Common Voice \parencite{anastassiou2024seedtts,ardila2020commonvoice}, under clean and noisy prompt conditions. Therefore, the findings may not apply to other languages or to very different input conditions.

The reliability guarantee holds only under the statistical assumptions of the calibration procedure. It limits the added risk of quality failures on average across inputs drawn from the same distribution as the calibration data, and it does not guarantee acceptable quality for every individual utterance. The guarantee may not hold when the inputs differ from the calibration data. In addition, running time is measured only on an NVIDIA Tesla T4 GPU through Google Colab, and no measurements are made on mobile or other limited-resource devices. Furthermore, the study evaluates only technical performance and computing efficiency. It does not assess clinical effectiveness, user experience, communication outcomes, or actual use in mobile assistive technologies. Listening tests with people are also outside the scope of the study, except for a possible small diagnostic check by the research team.

% =====================================================================
\subsection{Theoretical Framework}
\label{sec:theoretical}

This study rests on four bodies of theory, summarized in Figure~\ref{fig:theoretical}. Each one supplies one part of the proposed method, and together they explain why an earlier handover can be safe for some utterances and how that safety can be checked.

\paragraph{Flow matching} Flow matching treats generation as moving a sample of random noise toward realistic data by following a learned velocity field, from time $t=0$ (noise) to $t=1$ (data) \parencite{lipman2023flow,liu2023flow}. The models used in this study are trained on straight paths between noise and data. Because the paths are straight, one step of the remaining length along the current velocity gives an estimate of the final output from any point along the way, $\hat x_1(t)=x_t+(1-t)\,v(x_t,t)$, without any extra network evaluation \parencite{lipman2024guide}. F5-TTS, the teacher in this study, follows this approach \parencite{chen2024f5}. This theory gives the study two things: the endpoint estimate that the teacher hands to the student, and a way to observe, during generation, whether the teacher's prediction of the final speech has stopped changing. If it has, more teacher steps are expected to add little.

\paragraph{Knowledge distillation and teacher-guided sampling} Knowledge distillation trains a fast student model to reproduce the outputs of a slower teacher \parencite{yin2024dmd,li2024dmospeech}. A distilled student is cheap but tends to lose some variety in its outputs, which DMOSpeech~2 restores by letting the teacher run the first steps and the student finish \parencite{li2025dmospeech2}. The reasoning behind this division of labor is that the early, noisy steps decide the coarse structure of the speech, such as timing and prosody, while the later steps refine acoustic detail \parencite{chen2024f5,li2025dmospeech2}. At the handover, the student mixes the teacher's estimate with fresh noise before continuing; studies of this kind of re-noising show that it keeps the overall layout of the guiding input while correcting smaller errors \parencite{meng2022sdedit,xu2023restart}. Because each guided teacher step costs two network evaluations \parencite{ho2022cfg} and each student step costs one, the switching time largely decides the cost of an utterance. This theory explains why the handover point is where computation can be saved, and why the student can tolerate a teacher state that is not yet perfect.

\paragraph{Input-adaptive computation} Input-adaptive computation holds that inputs differ in difficulty, so a system can save computation by giving each input only as much processing as it needs and stopping once a signal shows that enough has been done \parencite{graves2016act,teerapittayanon2016branchynet}. Confident Adaptive Language Modeling applied this idea to language generation and used, among others, a signal that measures whether the model's internal state has stopped changing \parencite{schuster2022calm}. In this study, the same idea becomes a stopping rule for the teacher. At each candidate step, the rule checks the prompt SNR and the endpoint-stabilization score, and stops the teacher only when they indicate that an early handover is safe; otherwise the teacher continues to the reference step, so the rule never costs more than the fixed setting.

\paragraph{Statistical risk control} An adaptive rule will sometimes stop too early, so it needs a way to state how often this happens. Distribution-free risk control chooses the setting of a rule from held-out calibration data so that its expected rate of errors, the risk, stays below a chosen tolerance $\alpha$ with probability at least $1-\delta$ \parencite{bates2021rcps}. Learn-Then-Test does this by treating each candidate threshold as a hypothesis, ``the risk of this threshold is above $\alpha$'', and accepting only the thresholds for which the calibration data give strong evidence against that hypothesis \parencite{angelopoulos2025ltt}; testing the thresholds in a fixed order keeps the overall chance of error at $\delta$ \parencite{bauer1991}. The result is a guarantee of the form $\P(R\le\alpha)\ge1-\delta$, which holds on average over new inputs drawn like the calibration inputs, not for every single utterance. This theory supplies the calibration step and the reliability guarantee.

Taken together, these theories form the reasoning of the study. Flow matching makes the state of generation observable at every teacher step at no extra cost; teacher-guided sampling makes the handover point the main lever on cost; input-adaptive computation suggests setting that lever separately for each utterance; and statistical risk control makes it possible to state, before use, how often doing so will add a quality failure. The study tests whether this reasoning holds in practice for teacher-guided flow-matching TTS.

% Figure: theoretical framework (Chapter 1)
\begin{figure}[!htbp]
\centering
\begin{adjustbox}{max width=\linewidth}
\begin{tikzpicture}[
  font=\figfont,
  theory/.style={draw=frozenLine, fill=frozenFill, rounded corners=2pt, align=center, text width=3.7cm, minimum height=1.5cm, inner sep=3pt},
  gives/.style={draw=newLine, fill=newFill, rounded corners=2pt, align=center, text width=4.4cm, minimum height=1.5cm, inner sep=3pt},
  goal/.style={draw=newLine, very thick, fill=newFill, rounded corners=2pt, align=center, text width=3.2cm, minimum height=2.4cm, inner sep=4pt},
  head/.style={font=\figfont\bfseries, align=center},
  arr/.style={-{Stealth[length=2mm]}, thick}
]
\node[head] at (0,1.25) {Theory};
\node[head] at (5.0,1.25) {What it gives the study};
\node[head] at (10.0,1.25) {Proposed method};

\node[theory] (t1) at (0,0)     {\textbf{Flow matching}\\straight paths from noise to speech};
\node[theory] (t2) at (0,-2.0)  {\textbf{Knowledge distillation and teacher-guided sampling}};
\node[theory] (t3) at (0,-4.0)  {\textbf{Input-adaptive computation}\\stop when enough is done};
\node[theory] (t4) at (0,-6.0)  {\textbf{Statistical risk control}\\Learn-Then-Test};

\node[gives] (g1) at (5.0,0)    {Endpoint estimate at every teacher step, at no extra cost, and a score of whether it has settled};
\node[gives] (g2) at (5.0,-2.0) {A student that finishes from the teacher's state; the switching time sets the cost};
\node[gives] (g3) at (5.0,-4.0) {A stopping rule that gives each utterance only the teacher steps it needs};
\node[gives] (g4) at (5.0,-6.0) {A calibrated threshold with a bound on added failures: $\P(R\le\alpha)\ge1-\delta$};

\node[goal] (goal) at (10.0,-3.0) {\textbf{Risk-calibrated adaptive handover}\\[2pt]switching time chosen per utterance};

\foreach \i in {1,...,4} {
  \draw[arr] (t\i) -- (g\i);
  \draw[arr] (g\i.east) -- (goal.west);
}
\end{tikzpicture}
\end{adjustbox}
\caption{Theoretical framework of the study. Each body of theory supplies one part of the proposed method.}
\label{fig:theoretical}
\end{figure}


% =====================================================================
\subsection{Conceptual Framework}
\label{sec:conceptual}

The conceptual framework, shown in Figure~\ref{fig:conceptual}, follows an input--process--output model. Its process stage is organized by the four specific objectives, so each research question can be traced from what goes into the study to what comes out of it.

\paragraph{Input} The inputs are the test material, the models, and the settings fixed before any experiment. The test material consists of English zero-shot prompts and texts from the Seed-TTS test set, each under a clean and a noisy (15~dB) prompt condition, divided by speaker into development, calibration, and test sets. The models are the released F5-TTS teacher, the DMOSpeech~2 student, the duration predictor, and the vocoder, all with unchanged weights. The settings are the candidate switching times, the reference switching time, the rule that defines an added quality failure, and the tolerated risk $\alpha$ and confidence parameter $\delta$.

\paragraph{Process} The process follows the specific objectives. First, the study runs every candidate switching time on the development set and measures its quality and cost (Objective~1). Second, it states the choice of switching time as a cost-minimization problem with a limit on risk, and defines the stopping rule from the prompt SNR and the endpoint-stabilization score (Objective~2). Third, it calibrates the rule's threshold on the calibration set with Learn-Then-Test and checks the guarantee on the test set (Objective~3). Fourth, it compares the calibrated method with fixed switching times, an SNR-only rule, and an offline oracle on the test set, under clean and noisy prompts and with different calibration sizes (Objective~4).

\paragraph{Output} The outputs are a switching time chosen for each utterance and the speech it produces; a calibrated threshold with a stated bound on the rate of added quality failures; measured speech quality and computing cost for every compared method; and evidence on whether the endpoint-stabilization score adds value beyond the prompt SNR. These outputs answer the general problem: whether adaptive handover reduces computing cost while keeping added quality failures within the limit set beforehand.

% Figure: conceptual framework, input-process-output (Chapter 1)
\begin{figure}[!htbp]
\centering
\begin{adjustbox}{max width=\linewidth}
\begin{tikzpicture}[
  font=\figfont,
  io/.style={draw=frozenLine, fill=frozenFill, rounded corners=2pt, align=left, text width=3.9cm, minimum height=8.6cm, inner sep=5pt},
  outbox/.style={draw=newLine, fill=newFill, rounded corners=2pt, align=left, text width=3.9cm, minimum height=8.6cm, inner sep=5pt},
  ostep/.style={draw=newLine, very thick, fill=white, rounded corners=2pt, align=left, text width=4.3cm, minimum height=1.75cm, inner sep=4pt},
  head/.style={font=\figfont\bfseries, align=center},
  arr/.style={-{Stealth[length=2.2mm]}, very thick}
]
\node[head] at (0,1.0)   {INPUT};
\node[head] at (5.4,1.0) {PROCESS};
\node[head] at (10.8,1.0){OUTPUT};

\node[io, anchor=north] (in) at (0,0.6) {%
\textbf{Test material}\\
Seed-TTS English prompts and texts; clean and 15~dB noisy prompts; speaker-disjoint development, calibration and test sets\\[5pt]
\textbf{Models (weights unchanged)}\\
F5-TTS teacher, DMOSpeech~2 student, duration predictor, vocoder\\[5pt]
\textbf{Settings fixed in advance}\\
candidate and reference switching times; added-failure rule; tolerated risk $\alpha$ and confidence parameter $\delta$};

\node[ostep, anchor=north] (o1) at (5.4,0.6)   {\textbf{Objective 1.} Run every candidate switching time; measure quality and cost};
\node[ostep, below=0.45cm of o1] (o2) {\textbf{Objective 2.} State the cost--risk problem; define the stopping rule from prompt SNR and endpoint stabilization};
\node[ostep, below=0.45cm of o2] (o3) {\textbf{Objective 3.} Calibrate the threshold with Learn-Then-Test; check the guarantee on held-out speakers};
\node[ostep, below=0.45cm of o3] (o4) {\textbf{Objective 4.} Compare with fixed switching times, SNR-only rule and oracle; clean and noisy prompts; calibration sizes};

\node[outbox, anchor=north] (out) at (10.8,0.6) {%
\textbf{Per-utterance switching time} and the speech it produces\\[5pt]
\textbf{Calibrated threshold} with a stated bound on added quality failures\\[5pt]
\textbf{Quality and cost} of every compared method: WER, SIM, UTMOS, forward passes, runtime\\[5pt]
\textbf{Evidence} on whether endpoint stabilization adds value beyond prompt SNR};

\draw[arr] (o1) -- (o2);
\draw[arr] (o2) -- (o3);
\draw[arr] (o3) -- (o4);
\coordinate (pin)  at ($(o1.west)!0.5!(o4.west)$);
\coordinate (pout) at ($(o1.east)!0.5!(o4.east)$);
\draw[arr] (in.east |- pin) -- (pin);
\draw[arr] (pout) -- (out.west |- pout);
\end{tikzpicture}
\end{adjustbox}
\caption{Conceptual framework of the study in input--process--output form. The process stage follows the four specific objectives.}
\label{fig:conceptual}
\end{figure}


Table~\ref{tab:variables} lists the variables of the study. The switching time is the only part of the generation process that the method changes; everything else is held fixed, so that any difference in quality or cost between two runs of the same input can be traced to where the teacher stopped.

\begin{table}[!htbp]
\centering\small
\caption{Variables of the study.}
\label{tab:variables}
\begin{tabularx}{\textwidth}{@{}>{\raggedright\arraybackslash}p{2.6cm} >{\raggedright\arraybackslash}p{4.0cm} >{\raggedright\arraybackslash}X@{}}
\toprule
\textbf{Type} & \textbf{Variable} & \textbf{How it is set or measured} \\
\midrule
Independent & Switching time & Chosen for each utterance among four candidate teacher steps, the latest being the reference \\
Independent & Prompt condition & Clean, or with added noise at 15~dB SNR \\
Independent & Calibration sample size & 50, 100, or 200 speakers (Objective~4) \\
Decision signals & Prompt SNR; endpoint-stabilization score & Computed during generation with no extra network evaluation \\
Controlled & Model weights; student restart level; random seeds; hardware & Kept as released; default restart level; same seeds for every candidate of an input; Tesla T4 GPU \\
Dependent (quality) & WER, SIM, UTMOS, added-failure rate & Measured on every generated utterance; WER and SIM decide an added failure \\
Dependent (cost) & Forward passes; runtime and real-time factor (RTF) & Counted exactly per utterance; timed on the Tesla T4 GPU \\
\bottomrule
\end{tabularx}
\end{table}

% =====================================================================
\subsection{Definition of Terms}
\label{sec:terms}

The following terms are defined as they are used in this study.
\begin{description}[leftmargin=0pt,itemsep=2pt,font=\normalfont\itshape]
\item[Adaptive handover policy.] The rule that chooses a switching time for each utterance from the prompt SNR, the endpoint-stabilization score, or both, and uses the reference switching time when no earlier candidate qualifies.
\item[Added quality failure.] An utterance whose speech at a candidate switching time is clearly worse than at the reference switching time on the same input, judged by a word error rate and speaker similarity rule fixed before calibration.
\item[Calibration.] Choosing the policy's threshold on a separate set of utterances, not used for the final evaluation, so that the risk is bounded with a stated confidence.
\item[Candidate switching time.] One of the small, predefined set of switching times among which the policy chooses for each utterance.
\item[Endpoint estimate.] The teacher's prediction of the final mel-spectrogram, computed at any step from the current state and velocity without an extra network evaluation; it is what the teacher hands to the student.
\item[Endpoint-stabilization score.] How much the endpoint estimate changes between two consecutive teacher steps; a small score means the teacher's prediction of the final speech has settled.
\item[Flow matching.] A way of generating data in which a network learns the direction in which random noise should move, step by step, to become speech.
\item[Forward pass.] One evaluation of a neural network, also called a neural function evaluation. Computing cost is counted in forward passes, with a guided teacher step counting as two.
\item[Learn-Then-Test.] The statistical procedure used to choose the policy's threshold so that the risk stays below the tolerance with a stated confidence, under its assumptions.
\item[Oracle.] A reference rule, not usable in practice, that picks in hindsight the cheapest candidate switching time that does not fail for each utterance; it shows the largest possible saving.
\item[Prompt.] The short recording of a speaker given to the system, whose voice the generated speech should match.
\item[Prompt SNR.] An estimate, in decibels, of how loud the speech in the prompt is compared with its background noise.
\item[Reference switching time.] The fixed switching time of the baseline setting, against which the cost and quality of every other choice are compared.
\item[Reliability guarantee.] The statement produced by calibration that, with probability at least $1-\delta$, the risk of the chosen threshold is at most $\alpha$; it holds on average over inputs, not for each utterance.
\item[Restart level.] The point on the student's own schedule at which it starts generating from the teacher's estimate mixed with fresh noise; it is kept at the released default in this study.
\item[Risk.] The expected rate of added quality failures over new inputs.
\item[Speaker similarity.] How closely the voice in the generated speech matches the reference recording, measured by a speaker-verification model.
\item[Student model.] The smaller, distilled model that finishes generation in a few steps; in this study, the DMOSpeech~2 generator.
\item[Switching time.] The point in generation at which the teacher stops and the student takes over; also called the handover point.
\item[Teacher model.] The larger, multi-step model whose outputs the student was trained to imitate; in this study, F5-TTS.
\item[Teacher-guided sampling.] Generation in which the teacher performs the first steps and the student completes the output.
\item[Training-free.] Describes a method that changes only how the released models are run and does not train or fine-tune any network.
\item[UTMOS.] A model that predicts the mean opinion score listeners would give an utterance; it is used as an estimate of naturalness.
\item[Word error rate.] The share of words that a speech recognizer gets wrong when it transcribes the generated speech, compared with the input text; it measures intelligibility.
\item[Zero-shot TTS.] Speech generation in the voice of a speaker heard only through a short prompt, without training the model on that speaker.
\end{description}
\FloatBarrier